% !TEX root = userguide.tex

\def\headdate{10th October 2012}

\section{Eltree}
\label{sec:exampleeltree}

In this section examples are given for using the module eltree
(Sec.~\ref{sec:subdivide}).

The main purpose of the module-on eltree is to divide a reference domain/element
into subdomains to easily integrate on
parts of the domain and/or on the zero levelset interface.
A typical eltree generates a subdivision like in Figure~\ref{fig:eltree},
where subdomains are progressively smaller near a line or surface. Use is made of a levelset function to 
separate parts of the domain and refine near the interface between the parts.
\begin{figure}[hbtp!]
 \begin{center}
  \includegraphics[scale=0.5]{meshquadtree}
 \end{center}
  \caption{Subdivision near a circular interface using a quadtree of quadrilateral subdomains.} 
  \label{fig:eltree}
\end{figure}
The element tree can exist lines (1D), triangles, quads (2D), tetrahedrons or hexahedrons (3D).
The subdivision into triangles and hexahedrons is not straight forward and has been further
described in Appendix~\ref{chap:subdivide}.
If the domain is different from the reference element a mapping function
(\texttt{mapcoor}) needs to be defined.

The basic tree structure can be combined with a further subdivision of the elements
near the interface using the module \texttt{submesh} (see Sec.~\ref{sec:submesh}).
The module \texttt{intmesh} (see Sec.~\ref{sec:intmesh}) can optionally be used to
reconstruct the zero levelset interface and define a composite surface integration rule.
A typical eltree with further subdivision generates a subdivision like in Figure~\ref{fig:eltreesubmesh}.
\begin{figure}[hbtp!]
 \begin{center}
  \includegraphics[scale=0.5]{eltreesubmesh}
 \end{center}
  \caption{Subdivision near a circular interface using a quadtree and 
a further subdivision of the quads at the interface into triangles aligned with
the interface (zero levelset).} 
  \label{fig:eltreesubmesh}
\end{figure}
Note, that the subdivision is used for integration only and that therefore the
actual geometric properties of the triangles is unimportant as long as they fill the domain.

In the example file \texttt{eltree1.f90} an integral is performed over a circular (2D) or spherical (3D)
domain using a quadtree of quadrilaterals or octree of hexahedrons.

In the example file \texttt{eltree2.f90} a unit hexahedron is divided into a
number of tetrahedrons using the module submesh.
An integral is performed on part of the domain.

In the example file \texttt{eltree3.f90} an integral is performed over a circular (2D) or spherical (3D)
domain using a quadtree  of quadrilaterals or octree of hexahedrons with a further subdivision into triangles or tetrahedrons near the interface.

In the example file \texttt{eltree4.f90} an integral is performed over the boundary of a circular (2D) or spherical (3D)
domain using a quadtree of quadrilaterals  or octree of hexahedrons with a further subdivision into triangles or tetrahedrons near the interface.
The zero levelset interface is reconstructed and used for integration.
In Fig.~\ref{fig:sphere_surface} an example of a reconstructed surface mesh
of a sphere is shown.
\begin{figure}[hbtp!]
 \begin{center}
  \includegraphics[scale=0.3]{sphere_surface}
 \end{center}
  \caption{Reconstructed surface mesh of a sphere}
  \label{fig:sphere_surface}
\end{figure}

In the example file \texttt{eltree5.f90} an integral is performed over the boundary of a circular (2D)
domain using a finite element mesh, where each element is subdivided using a quadtree of quadrilaterals and further subdivided into triangles
near the interface. The zero levelset interface is reconstructed and used for integration (see Fig.~\ref{fig:mesh_sd5_plot} for
a rather coarse example).
\begin{figure}[hbtp!]
 \begin{center}
  \includegraphics[scale=0.5]{mesh_sd5_plot}
 \end{center}
  \caption{Curved $5\times5$ mesh, with a subdivision near a circular interface using a quadtree of quadrilaterals and 
a further subdivision of the quads at the interface into triangles aligned with
the interface (zero levelset). The red curve is the reconstructed zero levelset interface.} 
  \label{fig:mesh_sd5_plot}
\end{figure}

The example file \texttt{eltree6.f90} is similar to \texttt{eltree5.f90}, but now the problem is 3D and the integral in
on the surface of a sphere.

The examples \texttt{eltree7.f90}--\texttt{eltree10.f90} are similar to the examples above, but now using a quadtree
of triangles or an octree of tetrahedrons.
